	Desde a concepção da ideia de computadores, as pessoas têm se questionado sobre a possibilidade de essas máquinas se tornarem, um dia, inteligentes(Moons;Bankman;Verhelst, 2019). Atualmente, Inteligência artificial (IA) é um campo de pesquisa próspero com muitas aplicações práticas e tópicos de pesquisa.  Softwares inteligentes já são capazes de automatizar tarefas como compreensão e tradução de linguagens (Chiu et al. 2017), diagnósticos baseados em imagem na medicina (Esteva et al. 2017).
	
	Nos primórdios da inteligencia artificial, foram resolvidos rapidamente problemas que podem ser descritos como difíceis para seres humanos, mas que são relativamente fáceis e diretos para computadores, problemas que podem ser descritos por uma lista de regras matemáticas. Porém tarefas que são consideradas fáceis e intuitivas, para os humanos, porém difíceis de ser formalmente descritas, como reconhecer discurso ou imagens, são mais desafiadoras para as máquinas (Goodfellow et al. 2016).

O verdadeiro desafio se mostrou ser resolver tarefas que são consideradas fáceis, intuitivas, para os humanos e difíceis de ser formalmente descritas, como reconhecer discursos ou imagens 

	Nas últimas décadas, o campo de inteligência artificial fez avanços impressionantes na visão computacional com o rápido progresso na recognição de imagens.  Nos últimos anos o desempenho de algoritmos de detecção de objetos disparou de aproximadamente 30\% de precisão para mais de 90\% no benchmark PASCAL VOC . Para classificação de imagens no conjunto de dados ImageNet, os melhores algoritmos já são capazes de exceder o desempenho humano. Esse progresso no reconhecimento de imagens afeta várias aplicações como vigilância por vídeo, direção autônoma, assistência médica inteligente.(Jifeng Dai;Steve Lin ,  2018)
	Apesar do progresso da inteligência artificial na recognição de imagens, existem ainda grandes desafios a serem superados. Um desses desafios está em como treinar modelos que reconheçam bem as distintas configurações do mundo real que não foram vistas no treinamento da IA. Atualmente Classificadores Automáticos são treinados e avaliados com um conjunto de dados que é dividido aleatoriamente em conjuntos de treinamento e teste, assim as imagens no conjunto de teste possuem as mesmas condições  que o conjunto de treinamento. Porém, em aplicações no mundo real as imagens podem ter diferentes ângulos, escalas, configurações de câmera, etc. Essas diferenças no conjunto de dados podem levar a grandes quedas de precisão significantes em uma grande variedade de arquiteturas. Essa suscetibilidade dos modelos atuais pode causar erros severos em aplicações como a direção autônoma de veículos.
 	Este trabalho visa descobrir maneiras de reduzir o impacto da diversidade de condições das imagens na precisção dos modelos de classificadores automáticos, tornando-os mais adaptáveis à diversidade dos dados nos diferentes conjuntos de dados que podem ser encontrados nas mais variadas aplicações. Serão tomados como base os conceitos de IA e visão computacional encontrados na literatura para propor um método que torne os modelos atuais de classificação mais estáveis e confiáveis.
 	
 	Para que uma máquina possa aprender sobre o mundo real é necessária uma grande quantidade de dados para treinamento, de maneira que ele possa fazer generalizações a partir deles, e quando encontrar novos dados possa tomar decisões em cima dos mesmos. Assim como uma criança precisa interagir com o mundo ao seu redor para aprender como ele funciona , uma máquina precisa de uma enorme quantidade de dados sobre a aplicação para que ela possa entende-la, modela-la, e interagir com ela de acordo com o que aprendeu(Bengio et al, 2009). 

 	A inteligência artificial tem se tornado constantemente mais presente e relevante nos últimos anos. Diversos setores utilizam técnicas da inteligência computacional para interpretar dados e obter informações utéis a partir deles. Áreas como mercado financeiro,medicina, marketing digital, serviços de streaming, etc, se beneficiam de técnicas como o aprendizado de máquina para processar e interpretar dados e a partir deles fazer predições que auxiliem na tomada de decisões (GoodFellow et al., 2016). Softwares inteligentes já são capazes de automatizar tarefas como compreensão e tradução de linguagens (Chiu et al. 2017) e diagnósticos baseados em imagem na medicina (Esteva et al. 2017).

O aumento do volume de dados disponíveis em bases de dados específicas para aprendizado de máquina, a necessidade de extrair informações úteis de grandes quantidades de dados e, sobretudo, o aumento da capacidade de processamento dos computadores, resultante, principalmente, da evolução das arquiteturas de GPU e da programação paralela, incentivaram o desenvolvimento do modelo computacional de aprendizado de máquina conhecido como aprendizado profundo.(GoodFellow et al., 2016) 

O aprendizado em profundidade, ou deep learning, termo pelo qual é popularmente conhecido, é assim chamado pelo uso de  redes neurais multicamadas, responsáveis por tornar possível a extração de atributos de uma imagem, por exemplo, e a partir dessas características fazer predições. Com o crescimento dos banco de dados disponíveis para treinar modelos e os constantes avanços técnológicos na área de hardware e software, essa técnica passou a ser mais utilizadas em várias aplicações. (Goodfellow et al., 2016).

